% ============================================================
%  theorems.tex  —  Barcha teorema muhitlari
% ============================================================

% amsthm allaqachon packages.tex da yuklangan.

% ---  Uslublar ---
\newtheoremstyle{uzb_plain}%  nomi
  {8pt}%   yuqoridan bo'sh joy
  {6pt}%   pastdan bo'sh joy
  {\itshape}% matn fonti
  {}%      tekinlashtirish
  {\bfseries}% sarlavha fonti
  {.}%     sarlavhadan keyin belgisi
  { }%     sarlavha va matn orasidagi bo'sh joy
  {\thmname{#1}\thmnumber{ #2}\thmnote{ (#3)}}

\newtheoremstyle{uzb_definition}%  nomi
  {8pt}{}
  {\normalfont}
  {}
  {\bfseries}
  {.}{ }
  {\thmname{#1}\thmnumber{ #2}\thmnote{ (#3)}}

\newtheoremstyle{uzb_remark}%  nomi
  {6pt}{}
  {\normalfont}
  {}
  {\itshape}
  {.}{ }
  {\thmname{#1}\thmnumber{ #2}\thmnote{ (#3)}}

% --- Teorema muhitlari ---
\theoremstyle{uzb_definition}
\newtheorem{tarif}{Ta'rif}[section]
\newtheorem{aksioma}{Aksioma}[section]
\newtheorem{yasash}{Yasash}[section]   % Keyingi boblar uchun: geometrik yasash

\theoremstyle{uzb_plain}
\newtheorem{teorema}{Teorema}[section]
\newtheorem{lemma}{Lemma}[section]
\newtheorem{natija}{Natija}[section]

\theoremstyle{uzb_remark}
\newtheorem{izoh}{Izoh}[section]
\newtheorem{misol}{Misol}[section]     % Keyingi boblar uchun: sonli misollar

% --- tcolorbox ga o'ralgan muhitlar (ixtiyoriy foydalanish uchun) ---
\newenvironment{formulabox}[1][]{%
  \begin{tcolorbox}[formulabox, title={#1}]%
}{%
  \end{tcolorbox}%
}

\newenvironment{resultbox}[1][]{%
  \begin{tcolorbox}[resultbox, title={#1}]%
}{%
  \end{tcolorbox}%
}
